\section*{الفصل \ref{ch:g6:circuits-and-safety} --- الدارات الكهربائية والسلامة}

\begin{solution}{exo:g6:circuits-and-safety:1}
بأمانة: التوصيلات --- أيّ عنصر يتّصل بأيّ عنصر، وأين تتشعّب
الطرق. والمهمَل: أطوال الأسلاك الحقيقية وألوانها ومساراتها
المتلوّية على الطاولة.
\end{solution}

\begin{solution}{exo:g6:circuits-and-safety:2}
\omterm{def:g3:first-electric-circuit:battery}{البطارية}: خطّان متوازيان غير متساويين؛ \omterm{def:g3:first-electric-circuit:bulb}{والمصباح}: دائرة فيها
صليب؛ \omterm{ex:g3:first-electric-circuit:switch}{والقاطعة} المفتوحة: جسر مرفوع؛ \omterm{ex:g3:first-electric-circuit:switch}{والقاطعة} المغلقة: الجسر
موضوعًا مستويًا؛ \omterm{ex:g6:circuits-and-safety:motor}{والمحرّك}: حرف M في دائرة؛ \omterm{ex:g6:circuits-and-safety:motor}{والطنّان}: دائرة
مضروبة بخطّ. والخطّ \emph{الطويل} يدلّ على القطب $+$ في
\omterm{def:g3:first-electric-circuit:battery}{البطارية}.
\end{solution}

\begin{solution}{exo:g6:circuits-and-safety:3}
مستطيل واحد: \omterm{def:g3:first-electric-circuit:battery}{بطارية} وقاطعة ضغط وطنّان في حلقة واحدة. ورفع
الإصبع يفتح \omterm{ex:g3:first-electric-circuit:switch}{القاطعة} --- فيصير في الطريق الوحيد فراغ، ويُسكت
قانون الحلقة الطنّانَ في الحال.
\end{solution}

\begin{solution}{exo:g6:circuits-and-safety:4}
ليس بالضرورة. والسؤال الحاسم: مَن مُوصَّل بمَن؟ فالعناصر نفسها
بالتوصيلات نفسها --- الدارة نفسها، مهما اختلف منظر التشابك.
\end{solution}

\begin{solution}{exo:g6:circuits-and-safety:5}
تواصل المروحة دورانها: فموت \omterm{def:g3:first-electric-circuit:bulb}{المصباح} لا يُظلم إلا فرعه --- وهو
استقلال التفرع. أما \omterm{ex:g3:first-electric-circuit:switch}{القاطعة} المفتوحة فتجلس على الطريق المشترك
قبل التشعّب، فتوقف الاثنين: وهي قاعدة موقع الأمر.
\end{solution}

\begin{solution}{exo:g6:circuits-and-safety:6}
طريق \omterm{def:g4:conductors-insulators:conductor}{ناقل} عارٍ يصل جانبَي المنبع مباشرة متخطّيًا كل عنصر:
فيندفع التيار غير المعوَّق وتسخن الأسلاك. ومن ضحاياها: \omterm{def:g3:first-electric-circuit:battery}{البطارية}
(تُستنزف وتتلف)، وعزل السلك --- وعلى الشبكة أحيانًا البيت
كله.
\end{solution}

\begin{solution}{exo:g6:circuits-and-safety:7}
قطع فرعه في طرفة عين لأن التيار بدأ يندفع --- عطلٌ لا نزوة.
فجِد العطل وأصلِحه (انزع قابس المشتبه به، وافحص \omterm{def:g5:series-parallel-circuits:parallel}{الفرع}) قبل أن
تعيد الحارس.
\end{solution}

\begin{solution}{exo:g6:circuits-and-safety:8}
الحمّام: الجلد الرطب \omterm{def:g4:conductors-insulators:conductor}{ناقل}، ودفع الشبكة عبر جسم مبتلّ قد يكون
قاتلًا. والأسلاك المتضرّرة: العزل المتشقّق سياج فيه فجوة ---
فقد تلقى الأصابع النحاس. \omterm{def:g3:first-electric-circuit:battery}{والبطارية} نعم والمأخذ أبدًا: فدفع
\omterm{def:g3:first-electric-circuit:battery}{البطارية} ضئيل؛ ودفع المأخذ أشدّ منه بمئات المرات.
\end{solution}

\begin{solution}{exo:g6:circuits-and-safety:9}
المخطط: \omterm{def:g3:first-electric-circuit:battery}{بطارية}، وقاطعة مغلقة على الطريق المشترك، ثم تشعّب ---
\omterm{def:g5:series-parallel-circuits:parallel}{فرع} \omterm{ex:g6:circuits-and-safety:motor}{الطنّان} \omterm{def:g5:series-parallel-circuits:parallel}{وفرع} \omterm{def:g3:first-electric-circuit:bulb}{المصباح} --- يلتقيان عائدين إلى $-$. والمتوقّع:
أن يسطع \omterm{def:g3:first-electric-circuit:bulb}{المصباح} \emph{ويصوّت} \omterm{ex:g6:circuits-and-safety:motor}{الطنّان}، مستقلّين. ومع قاطعة
إضافية مفتوحة على \omterm{def:g5:series-parallel-circuits:parallel}{فرع} \omterm{ex:g6:circuits-and-safety:motor}{الطنّان}: صمت هناك، ويظلّ \omterm{def:g3:first-electric-circuit:bulb}{المصباح} ساطعًا
--- فقواطع الفروع لا تأمر إلا فروعها.
\end{solution}

\begin{solution}{exo:g6:circuits-and-safety:10}
السلك الإضافي \omterm{def:g6:circuits-and-safety:short}{دارة قصيرة}: طريق عارٍ مباشر من $+$ إلى $-$.
فيفضّل التيار الطريق السهل الخالي على \omterm{ex:g6:circuits-and-safety:motor}{المحرّك} الكادح، فتجوع
المروحة بينما يحمل سلك الاختصار الاندفاعَ ويسخن --- فيستنزف
\omterm{def:g3:first-electric-circuit:battery}{البطارية} ويدعو إلى الحروق. فانزع سلك ``المتانة''.
\end{solution}

\begin{solution}{exo:g6:circuits-and-safety:11}
\omterm{ex:g3:first-electric-circuit:switch}{القاطعة} الجدارية تفتح حلقة التجهيزة نفسها --- فيُقطع طريق
\omterm{def:g3:first-electric-circuit:bulb}{المصباح}. أما قاطع \omterm{def:g5:series-parallel-circuits:parallel}{الفرع} فيفتح \omterm{def:g5:series-parallel-circuits:parallel}{الفرع} كله أبعد في المنبع ---
فحتى لو كانت أسلاك التجهيزة معطوبة أو كانت \omterm{ex:g3:first-electric-circuit:switch}{القاطعة} الجدارية
مخطئة اللافتة، يبقى \omterm{def:g5:series-parallel-circuits:parallel}{الفرع} ميتًا. فراغان \omterm{def:g4:series-circuits:series}{على التسلسل}: كل واحد
وحده يوقف التيار، والاثنان معًا يغفران خطأً في أيّ منهما.
\end{solution}

\begin{solution}{exo:g6:circuits-and-safety:12}
\omterm{def:g3:first-electric-circuit:battery}{بطارية}، ثم \omterm{ex:g3:first-electric-circuit:switch}{القاطعة} الرئيسية على الطريق المشترك، ثم أربعة فروع
متفرّعة: \omterm{def:g3:first-electric-circuit:bulb}{مصباح} الخشبة 1؛ \omterm{def:g3:first-electric-circuit:bulb}{ومصباح} الخشبة 2؛ ومحرّك الستارة مع
قاطعته \omterm{def:g4:series-circuits:series}{على التسلسل}؛ \omterm{ex:g6:circuits-and-safety:motor}{والطنّان} مع قاطعته \omterm{def:g4:series-circuits:series}{على التسلسل}. والتحقّقات:
الرئيسية تأمر الجميع (الطريق المشترك)؛ والمصباحان يتعطّلان
مستقلّين (فرعان منفصلان)؛ \omterm{ex:g6:circuits-and-safety:motor}{والمحرّك} \omterm{ex:g6:circuits-and-safety:motor}{والطنّان} يطيع كل واحد
قاطعته (تسلسل داخل فرعه) ويتجاهل الآخر (تفرع بين الفروع).
\end{solution}

\begin{solution}{pb:g6:circuits-and-safety:1}
\textbf{1.} \omterm{def:g4:series-circuits:series}{على التسلسل} كان عنصر واحد ميت أو مطفأ يُظلم كل
شيء، وكانت العناصر الثلاثة تتقاسم الدفع فتخفت وتضعف كلها.
والتفرع يُبقي كل عنصر مستقلًّا وبكامل قوّته.
\textbf{2.} كل قاطعة تأمر عنصر فرعها هو بالضبط؛ ولا تأمر أيّ
واحدة منها الفروع الأخرى.
\textbf{3.} على الطريق المشترك، بين \omterm{def:g3:first-electric-circuit:battery}{البطارية} وأول تشعّب ---
قبل كل \omterm{def:g5:series-parallel-circuits:parallel}{فرع}.
\textbf{4.} \omterm{def:g3:first-electric-circuit:battery}{بطارية} $\to$ قاطعة رئيسية $\to$ \omterm{def:g5:series-parallel-circuits:parallel}{عقدة} تتشعّب إلى
ثلاثة فروع --- (\omterm{def:g3:first-electric-circuit:bulb}{مصباح} السقف $+$ قاطعته)، و(\omterm{def:g3:first-electric-circuit:bulb}{مصباح} الطاولة $+$
قاطعته)، و(محرّك المروحة $+$ قاطعته) --- ثم تلتقي عائدة إلى
\omterm{def:g3:first-electric-circuit:battery}{البطارية}.
\textbf{5.} هما \omterm{def:g4:series-circuits:series}{على التسلسل} في \omterm{def:g5:series-parallel-circuits:parallel}{فرع} واحد: فيتقاسمان الدفع
(فيضعفان معًا --- وخفوت \omterm{def:g3:first-electric-circuit:bulb}{مصباح} السقف كان في الحقيقة عرضًا من
أعراض هذا الزوج)، وقاطعة المروحة، الجالسة على طريقهما المشترك،
تأمرهما معًا.
\textbf{6.} افصِل بينهما: مُدّ سلكًا جديدًا واحدًا حتى ينال
\omterm{def:g3:first-electric-circuit:bulb}{مصباح} الطاولة فرعه الخاص عبر المغذّي --- فيصير كل واحد بكامل
قوّته ومستقلًّا.
\textbf{7.} \omterm{def:g6:circuits-and-safety:short}{دارة قصيرة} عبر القطبين. وافحص هل ما زالت \omterm{def:g3:first-electric-circuit:battery}{البطارية}
تعطي --- فالدارات القصيرة تستنزف وقد تُتلف الخلايا (\omterm{def:g3:first-electric-circuit:battery}{وبطارية}
سخنت تستحقّ التقاعد).
\textbf{8.} الدارات القصيرة (أسلاك عارية يلقى بعضها بعضًا)
والصعقات أو الحروق (أصابع تلقى نحاسًا عاريًا).
\textbf{9.} الجلد الرطب ينقل خيرًا بكثير من الجافّ؛ وحتى بناء
\omterm{def:g3:first-electric-circuit:battery}{البطارية} يستحقّ انضباط الأيدي الجافّة --- وفي البيت تلقى
اللمسة نفسها دفع الشبكة، الأشدّ بمئات المرات: وقد يكون قاتلًا.
فالعادات تُدرَّب على الدارة الآمنة.
\textbf{10.} كشف القاطع أن تيارًا مفرطًا كان يندفع في فرعه ---
عطلًا، على الأرجح مطرًا في تجهيزة. والترتيب الحكيم: اترك
القاطع مطفأً، وانزع القوابس وافحص \omterm{def:g5:series-parallel-circuits:parallel}{الفرع}، وأصلِح العطل أو جفّفه،
ثم أعِده بعد ذلك فقط.
\textbf{11.} حلقته الخاصة، منفصلة عن شبكة الكوخ (أو فرعًا
متفرّعًا آخر على \omterm{def:g3:first-electric-circuit:battery}{بطارية} الكوخ): \omterm{def:g3:first-electric-circuit:battery}{بطارية}، وقاطعة ضغط عند البيت،
وسلك مزدوج طويل، وطنّان عند الكوخ، ثم عودة --- فجرس الباب حلقة
تسلسل خاصة به مهما امتدّت أسلاكه.
\textbf{12.} مثلًا: ``المآخذ: تجاربنا تنتهي عند \omterm{def:g3:first-electric-circuit:battery}{البطارية} ---
فالمأخذ ليس لنا. والماء: الأيدي المبتلّة لا تلمس أسلاكًا، في
الداخل ولا في الخارج. والأسلاك العارية: كل وصلة ملفوفة؛ والسلك
العاري عطل لا اختصار.''
\end{solution}
